\section{Anlageurteil}
\label{sec:urteil}

Dieser Teil f\"uhrt die Befunde zusammen. Er f\"uhrt keine neue Zahl ein.

\subsection{Was belegt ist}

\begin{enumerate}
\item \textbf{Die Ertragslage ist die beste der Unternehmensgeschichte.}
Umsatz \MioUSDexakt{\Umsatz}, bereinigtes Ergebnis
\MioUSDexakt{\BereinigtesErgebnis}, freier Cash-Flow
\MioUSDexakt{\FCFReiheXXV} -- jeweils H\"ochstwerte der Zehnjahresreihe
(Tabelle~\ref{tab:reihe}).

\item \textbf{Sie stammt \"uberwiegend aus dem Goldpreis.} Die Produktion ging
um \PctBetrag{\ProduktionDelta} zur\"uck, die St\"uckkosten stiegen um
\Pct{\AISCDelta}, der erzielte Preis stieg um \Pct{\GoldpreisDelta}
(Abschnitt~\ref{sec:margejeunze}).

\item \textbf{Die Bilanz ist unbelastet.} Nettoliquidit\"at
\MioUSDexakt{\Nettoliquiditaet}, Eigenkapitalquote
\Pct{\Eigenkapitalquote}, unbeanspruchte Kreditlinie
\MioUSDexakt{\FazilitaetUnbeansprucht} (Abschnitt~\ref{sec:bilanz}).

\item \textbf{Das Unternehmen w\"achst aus eigener Substanz.} Reserven
\MozReserven{\Reserven} nach \Pct{\ReservenDelta} Zuwachs, im
\num{\ReservenJahreFolge}.\ Jahr in Folge steigend; ver\"offentlichter
Dreijahresausblick \Pct{\ProgAchtProdDelta} mehr Produktion und
\Pct{\ProgAchtAISCRueckgang} niedrigere St\"uckkosten bis 2028
(Teil~\ref{sec:strategie}).

\item \textbf{Der freie Cash-Flow der Reihe hat keinen Boden unter sich.}
Median \MioUSDexakt{\FCFMedian} \"uber zehn Jahre, kleinster Wert
\MioUSDexakt{\FCFMin}; erst die letzten drei Jahre sind deutlich positiv
(Abschnitt~\ref{sec:reihe}).

\item \textbf{Der Einbruch 2026 war \"uberwiegend kein Alamos-Ereignis, aber
nicht nur ein Marktereignis.} Aktie \Pct{\EinbruchAktie}, Sektor
\Pct{\EinbruchSektor}, Gold \Pct{\EinbruchGold}; der Abstand zum Sektor ist
mit \Pct{\AbstandUeberschuss} bis heute nicht aufgeholt
(Abschnitt~\ref{sec:zerlegung}).

\item \textbf{Keiner der gerechneten F\"alle erreicht den eigenen
Ma{\ss}stab.} Interner Zinsfu{\ss} zwischen
\Pct{\FlowIRRAFuenf} und \Pct{\FlowIRRCFuenfzehn} gegen eine H\"urde von
\Pct{\FlowHuerde} (Tabelle~\ref{tab:flowirr}).
\end{enumerate}

\subsection{Was das Urteil tr\"agt}

\bk{Es gibt in diesem Fall keinen Streit \"uber die Qualit\"at des
Unternehmens. Alamos f\"ordert Gold aus wachsenden Reserven, h\"alt mehr Kasse
als Schulden und hat einen Ausbauplan, den es aus dem laufenden Gesch\"aft
bezahlen kann, solange der Goldpreis nicht auf den Stand von 2024
zur\"uckf\"allt. Was zur Debatte steht, ist allein der Preis.

Der Streitpunkt l\"asst sich auf einen Satz bringen. Die Rechnung dieses
Dokuments schreibt einen Goldpreis konstant fort und rechnet den Ausbau, den
das Unternehmen ver\"offentlicht hat, mit hinein; sie kommt zu dem Schluss,
dass der heutige Kurs auch im g\"unstigsten belegbaren Fall zu hoch ist. Wer
dagegen h\"alt, muss einen h\"oheren Goldpreis oder einen h\"oheren Endwert
unterstellen als jede Zahl, die in den Prim\"arquellen steht. Das ist eine
zul\"assige Erwartung, aber es ist eine Erwartung und kein Befund.}

\subsection{Nicht verf\"ugbare Angaben}
\label{sec:nichtverfuegbar}

Diese Angaben fehlen. Sie werden benannt und nicht gesch\"atzt.

\begin{description}
\item[Die k\"unftige Entwicklung des Goldpreises.] Sie ist die tragende
Gr\"o{\ss}e dieses Falls und steht in keiner Prim\"arquelle. Deshalb rechnet
Teil~\ref{sec:flow} drei vom Unternehmen selbst erzielte Preise durch, statt
einen vierten zu erfinden.

\item[Die Kapitalkosten des Unternehmens.] Beta, risikoloser Zins und
Marktrisikopr\"amie stehen in keinem der ausgewerteten Dokumente. Ein daraus
konstruierter Kapitalkostensatz w\"are eine Erfindung mit dem Anschein einer
Berechnung. An seine Stelle tritt die belegte Eigenkapitalrendite des
Unternehmens selbst (Abschnitt~\ref{sec:anleger}).

\item[Die einzelnen Sch\"atzungen der Analysten.] Verf\"ugbar sind das
aggregierte Votum, das mittlere Kursziel und dessen Spanne
(Abschnitt~\ref{sec:konsens}); die zugrunde liegenden Modelle, die
Goldpreisannahmen und die Zuordnung der Ziele zu einzelnen H\"ausern sind es
nicht. Der Vergleich in Abschnitt~\ref{sec:begruendung} beschr\"ankt sich
deshalb auf Richtung und Niveau.

\item[Die Aufteilung des Streubesitzes.] Siehe Abschnitt~\ref{sec:luecken}.

\item[Der Investorentag vom \Datum{4.\,Februar 2026}.] Die Pr\"asentationen
enthalten nach Angabe des Unternehmens die ausf\"uhrliche Fassung des
Dreijahresausblicks.\quelleMDA{2025}{12}\quelleMerken{s20mdaxcacfxbc} Ausgewertet wurde nur, was
im Lagebericht selbst steht; die Pr\"asentation ist keine bei der
Wertpapieraufsicht eingereichte Unterlage und tr\"agt keine gedruckte Seite.
Das ist eine bewusste Beschr\"ankung und keine \"Ubersehung.
\end{description}
